\section{Encontro, aproximação e estimação de pose}

As operações de rendezvous em ambiente orbital envolvem desafios que ultrapassam o simples problema de deslocamento relativo entre duas espaçonaves. Em missões de serviço em órbita, inspeção, captura e berthing, a operação depende da articulação entre navegação relativa, definição de trajetórias seguras, estimação da pose do alvo e controle da aproximação final. Esses elementos tornam-se fortemente interdependentes, sobretudo em cenários de maior proximidade, nos quais erros de percepção, navegação ou controle podem comprometer a segurança e o sucesso da missão \cite{arantes2010far,fiori2024modeling}.

Entre os trabalhos relevantes nessa linha, destaca-se o estudo de Arantes et al., que analisa manobras far e proximity para uma constelação de satélites de serviço e incorpora, no mesmo contexto, estratégias de transferência orbital e estimação autônoma da pose de um satélite cliente não cooperativo \cite{arantes2010far}. Nesse trabalho, as manobras distantes são tratadas com base na formulação de Lambert, enquanto a fase de aproximação e trajetória final é analisada por meio das equações de Hill. Além disso, durante a fase de inspeção, a pose relativa do alvo é estimada por visão computacional monocular, utilizando informações geométricas previamente conhecidas do satélite cliente \cite{arantes2010far}.

Esse resultado evidencia que, em operações de rendezvous, o problema não se limita ao planejamento da trajetória relativa. Em alvos não cooperativos, a aproximação segura exige também a capacidade de reconstruir, com precisão adequada, a posição e a orientação relativas do alvo ao longo da missão. No trabalho de Arantes et al., a estimação de pose é formulada como um problema de minimização associado à correspondência entre características visuais do alvo e sua projeção no sistema de imagem, permitindo integrar percepção visual e navegação relativa em um mesmo arcabouço operacional \cite{arantes2010far}.

Além da vertente associada à percepção, a literatura também trata o rendezvous sob a ótica do guiamento e do controle em proximidade. O trabalho de Fiori et al. propõe um arcabouço de modelagem, simulação e controle para o movimento de uma pequena espaçonave nas proximidades de uma estação espacial, considerando restrições posicionais e a presença de obstáculos físicos \cite{fiori2024modeling}. A abordagem emprega potenciais virtuais e um formalismo geométrico baseado em grupos de Lie para descrever tanto o movimento roto-translacional quanto os campos de controle. Embora esse estudo não trate diretamente de manipuladores robóticos, ele contribui para a compreensão de estratégias de aproximação final em cenários nos quais a convergência para o alvo deve ser compatibilizada com restrições geométricas e operacionais \cite{fiori2024modeling}.

De forma mais ampla, a literatura mostra que o problema de rendezvous pode ser tratado em diferentes níveis. Em alguns trabalhos, a ênfase recai sobre as manobras orbitais e a dinâmica relativa; em outros, sobre a navegação local, a percepção do alvo ou o controle em regiões de maior sensibilidade operacional. Observa-se, assim, um avanço gradual na integração entre trajetória relativa, percepção e guiamento autônomo, embora muitos estudos ainda privilegiem apenas uma dessas dimensões \cite{arantes2010far,fiori2024modeling}.

Essa observação é relevante para a presente dissertação, pois reforça a necessidade de tratar de forma articulada a dinâmica relativa entre perseguidor e alvo, a fase de aproximação curta e os requisitos associados à interação final em operações de berthing. Em particular, interessa aqui a formulação de um problema em que a aproximação orbital, a atitude da espaçonave perseguidora e a presença do manipulador robótico sejam considerados de maneira consistente dentro do mesmo contexto de missão.